\chapter{Observable and Mechanism Mapping}

This appendix records how measured or simulated observables constrain the
COMSOL model states. The mapping is not one-to-one; it identifies which
measurements best constrain each physical mechanism.

\begin{table}[H]
\centering
\begin{tabular}{p{0.22\textwidth}p{0.34\textwidth}p{0.32\textwidth}}
\toprule
Observable & Primary sensitivity & COMSOL use \\
\midrule
Normalized PCE & aggregate performance loss & aging trajectory and $T_{80}$ marker \\
$V_{oc}$ & recombination and electrostatic shifts & trap density, surface recombination, and charge calibration \\
$J_{sc}$ & collection and absorber-quality changes & photogeneration, collection, and strong recombination check \\
FF & leakage, shunting, contact, and curve-shape losses & $R_s$, $R_{sh}$, and interface-loss calibration \\
$V_{mpp}$, $J_{mpp}$ & operating-point shift & maximum-power-point and deployment relevance \\
Hysteresis index & mobile ions and field redistribution & $c_{\mathrm{ion},0}$, $\mu_{\mathrm{ion}}$, and boundary-condition calibration \\
PL peak shift & composition and bandgap change & absorber-state and photogeneration-loss constraint \\
PL spread & spatial heterogeneity & check against the one-dimensional uniformity assumption \\
XRD phase fraction & composition or secondary phases & absorber degradation and phase-segregation constraint \\
Impedance or transient current & ions, capacitance, and interface dynamics & ion mobility, ion density, and interface-state calibration \\
\bottomrule
\end{tabular}
\caption{Observation mapping for COMSOL degradation calibration.}
\label{tab:feature_state_mapping}
\end{table}

The table defines calibration targets rather than fitted proof. Terminal J--V
data are necessary but not sufficient because several mechanisms can produce
similar PCE loss. Hysteresis, impedance, transient current, PL, and XRD
measurements provide additional constraints that help separate ionic,
recombination, contact, leakage, and absorber-quality changes.
